% GENERATED FILE. Edit canonical lessons, then run npm run generate:books.
% canonical-source-hash: fnv1a64:db4851012dd2c198
% canonical-lessons: PA-C215-ask-return, PA-C215-ki-pinde, PA-C215-ki-review, PA-C215-exchange, PA-C215-payoff

\chapter{Do You Drink Tea?}
\label{ch:pa-a2-polar-habitual}
\hlchaptermodality{215}

\begin{chapteropening}
\textbf{By the end of this chapter:} \emph{I can ask whether someone does something every day and answer yes or no in a short exchange.}

Ask whether someone drinks, eats or goes somewhere, and answer with a whole yes or no sentence.
\end{chapteropening}

\section[\texorpdfstring{\pa{ਕੀ} \pa{ਤੁਸੀਂ} \pa{ਯਾਤਰੀ} \pa{ਹੋ}? … (kī tusīṁ yātrī …)}{ਕੀ ਤੁਸੀਂ ਯਾਤਰੀ ਹੋ? … (kī tusīṁ yātrī …)}]{Return to asking and saying no}

\label{lesson:PA-C215-ask-return}



\begin{quote}
\emph{Hear:} \emph{are you a traveller?} Say \textbf{\pa{ਕੀ} \pa{ਤੁਸੀਂ} \pa{ਯਾਤਰੀ} \pa{ਹੋ}?} Then answer: \textbf{\pa{ਹਾਂ} \pa{ਜੀ}, \pa{ਮੈਂ} \pa{ਯਾਤਰੀ} \pa{ਹਾਂ।}}
\end{quote}

\begin{grammarlens}[title={Grammar Lens}]
You can also say \textbf{\pa{ਮੈਂ} \pa{ਚਾਹ} \pa{ਨਹੀਂ} \pa{ਪੀਂਦਾ} \pa{ਹਾਂ}} and \textbf{\pa{ਤੁਸੀਂ} \pa{ਚਾਹ} \pa{ਪੀਂਦੇ} \pa{ਹੋ}}. The next lesson puts \textbf{\pa{ਕੀ}} in front of the second one.
\end{grammarlens}

\subsection*{Your turn}
\begin{itemize}
\raggedright
  \item \emph{Hear:} \emph{is the garden open?} Say \textbf{\pa{ਕੀ} \pa{ਬਾਗ਼} \pa{ਖੁੱਲ੍ਹਾ} \pa{ਹੈ}?}
  \item \emph{Hear:} \emph{I don't drink tea} Say it with the ending that fits you.
  \item \emph{Hear:} \emph{you drink tea} Say \textbf{\pa{ਤੁਸੀਂ} \pa{ਚਾਹ} \pa{ਪੀਂਦੇ} \pa{ਹੋ।}}
\end{itemize}

\begin{tcolorbox}[breakable,colback=teal!4,colframe=teal!35!black,title={Before you move on}]
Ask one yes-or-no question, then say one thing you don't do.
\end{tcolorbox}
\section[\texorpdfstring{\pa{ਕੀ} \pa{ਤੁਸੀਂ} \pa{ਚਾਹ} \pa{ਪੀਂਦੇ} \pa{ਹੋ}? (kī tusīṁ chāh …)}{ਕੀ ਤੁਸੀਂ ਚਾਹ ਪੀਂਦੇ ਹੋ? (kī tusīṁ chāh …)}]{\pa{ਕੀ} \pa{ਤੁਸੀਂ} \pa{ਚਾਹ} \pa{ਪੀਂਦੇ} \pa{ਹੋ}? — do you drink tea?}

\label{lesson:PA-C215-ki-pinde}



\begin{quote}
Say \emph{you drink tea}. Then \emph{are you a traveller?}
\end{quote}

\begin{grammarlens}[title={Grammar Lens: \pa{ਕੀ} and the every-day present}]
Nothing new to build: put \textbf{\pa{ਕੀ}} in front of a present sentence, and let the voice rise.

\begin{quote}
\textbf{\pa{ਕੀ} \pa{ਤੁਸੀਂ} \pa{ਚਾਹ} \pa{ਪੀਂਦੇ} \pa{ਹੋ}?} — \emph{kī tusīṁ chāh pīnde ho?} — Do you drink tea?
\end{quote}

\begin{quote}
\textbf{\pa{ਕੀ} \pa{ਉਹ} \pa{ਰੋਜ਼} \pa{ਸਕੂਲ} \pa{ਜਾਂਦੀ} \pa{ਹੈ}?} — \emph{kī uh roz sakūl jāndī hai?} — Does she go to school every day?
\end{quote}

\begin{quote}
\textbf{\pa{ਕੀ} \pa{ਤੁਸੀਂ} \pa{ਰੋਜ਼} \pa{ਕੰਮ} \pa{ਕਰਦੇ} \pa{ਹੋ}?} — \emph{kī tusīṁ roz kamm karde ho?} — Do you work every day?
\end{quote}
\end{grammarlens}

\subsection*{Your turn}
\begin{itemize}
\raggedright
  \item \emph{Hear:} \emph{do you eat roti?} Say \textbf{\pa{ਕੀ} \pa{ਤੁਸੀਂ} \pa{ਰੋਟੀ} \pa{ਖਾਂਦੇ} \pa{ਹੋ}?}
  \item \emph{Hear:} \emph{does he drink milk?} Say \textbf{\pa{ਕੀ} \pa{ਉਹ} \pa{ਦੁੱਧ} \pa{ਪੀਂਦਾ} \pa{ਹੈ}?}
  \item \emph{Read it:} \textbf{\pa{ਕੀ} \pa{ਤੁਸੀਂ} \pa{ਚਾਹ} \pa{ਪੀਂਦੇ} \pa{ਹੋ}?} Give its meaning.
  \item \emph{Write it:} copy \textbf{\pa{ਕੀ} \pa{ਤੁਸੀਂ} \pa{ਚਾਹ} \pa{ਪੀਂਦੇ} \pa{ਹੋ}?} once
\end{itemize}

\begin{tcolorbox}[breakable,colback=teal!4,colframe=teal!35!black,title={Before you move on}]
Ask someone whether they drink tea, and whether they go to school every day.
\end{tcolorbox}
\section[\texorpdfstring{\pa{ਕੀ} \pa{ਤੁਸੀਂ} \pa{ਚਾਹ} \pa{ਪੀਂਦੇ} \pa{ਹੋ}? (kī tusīṁ chāh …)}{ਕੀ ਤੁਸੀਂ ਚਾਹ ਪੀਂਦੇ ਹੋ? (kī tusīṁ chāh …)}]{Ask do-you questions again}

\label{lesson:PA-C215-ki-review}



\begin{quote}
\emph{Hear:} \emph{do you drink tea?} Say \textbf{\pa{ਕੀ} \pa{ਤੁਸੀਂ} \pa{ਚਾਹ} \pa{ਪੀਂਦੇ} \pa{ਹੋ}?}
\end{quote}

\subsection*{Your turn}
\begin{itemize}
\raggedright
  \item \emph{Hear:} \emph{does she go to school every day?} Say \textbf{\pa{ਕੀ} \pa{ਉਹ} \pa{ਰੋਜ਼} \pa{ਸਕੂਲ} \pa{ਜਾਂਦੀ} \pa{ਹੈ}?}
  \item \emph{Hear:} \emph{yes, she goes to school every day} Say \textbf{\pa{ਹਾਂ} \pa{ਜੀ}, \pa{ਉਹ} \pa{ਰੋਜ਼} \pa{ਸਕੂਲ} \pa{ਜਾਂਦੀ} \pa{ਹੈ।}}
  \item \emph{Read it:} \textbf{\pa{ਕੀ} \pa{ਤੁਸੀਂ} \pa{ਰੋਟੀ} \pa{ਖਾਂਦੇ} \pa{ਹੋ}?} Give its meaning.
  \item \emph{Write it:} cover the model and write \textbf{\pa{ਕੀ} \pa{ਉਹ} \pa{ਦੁੱਧ} \pa{ਪੀਂਦਾ} \pa{ਹੈ}?} once
\end{itemize}

\begin{tcolorbox}[breakable,colback=teal!4,colframe=teal!35!black,title={Before you move on}]
Ask one do-you question and answer it yourself.
\end{tcolorbox}
\section[\texorpdfstring{\pa{ਨਹੀਂ} \pa{ਜੀ}, \pa{ਮੈਂ} \pa{ਚਾਹ} … (nahīṁ jī, maiṁ …)}{ਨਹੀਂ ਜੀ, ਮੈਂ ਚਾਹ … (nahīṁ jī, maiṁ …)}]{A short exchange: tea or milk?}

\label{lesson:PA-C215-exchange}



\begin{quote}
Ask \emph{do you drink tea?} Answer it \emph{no}, in a whole sentence.
\end{quote}

\begin{grammarlens}[title={Grammar Lens: four lines}]
A host and a guest. Hear the turns; each answer repeats the verb.

\begin{quote}
\textbf{\pa{ਕੀ} \pa{ਤੁਸੀਂ} \pa{ਚਾਹ} \pa{ਪੀਂਦੇ} \pa{ਹੋ}?} — \emph{kī tusīṁ chāh pīnde ho?} — Do you drink tea?
\end{quote}

\begin{quote}
\textbf{\pa{ਨਹੀਂ} \pa{ਜੀ}, \pa{ਮੈਂ} \pa{ਚਾਹ} \pa{ਨਹੀਂ} \pa{ਪੀਂਦੀ} \pa{ਹਾਂ।}} — \emph{nahīṁ jī, maiṁ chāh nahīṁ pīndī hāṁ.} — No, I don't drink tea.
\end{quote}

\begin{quote}
\textbf{\pa{ਕੀ} \pa{ਤੁਸੀਂ} \pa{ਦੁੱਧ} \pa{ਪੀਂਦੇ} \pa{ਹੋ}?} — \emph{kī tusīṁ duddh pīnde ho?} — Do you drink milk?
\end{quote}

\begin{quote}
\textbf{\pa{ਹਾਂ} \pa{ਜੀ}, \pa{ਮੈਂ} \pa{ਦੁੱਧ} \pa{ਪੀਂਦੀ} \pa{ਹਾਂ।}} — \emph{hāṁ jī, maiṁ duddh pīndī hāṁ.} — Yes, I drink milk.
\end{quote}
\end{grammarlens}

\subsection*{Your turn}
\begin{itemize}
\raggedright
  \item \emph{Hear:} \emph{do you drink tea?} Answer \emph{no} as yourself.
  \item \emph{Hear:} \emph{do you drink milk?} Answer \emph{yes} as yourself.
  \item \emph{Read it:} the four lines Say both parts aloud.
  \item \emph{Write it:} copy the second line once
\end{itemize}

\begin{tcolorbox}[breakable,colback=teal!4,colframe=teal!35!black,title={Before you move on}]
Play the host: ask both questions. Then play the guest.
\end{tcolorbox}
\section[\texorpdfstring{\pa{ਕੀ} \pa{ਤੁਸੀਂ} \pa{ਚਾਹ} \pa{ਪੀਂਦੇ} \pa{ਹੋ}? (kī tusīṁ chāh …)}{ਕੀ ਤੁਸੀਂ ਚਾਹ ਪੀਂਦੇ ਹੋ? (kī tusīṁ chāh …)}]{Tea or milk: a four-skill payoff}

\label{lesson:PA-C215-payoff}



\begin{quote}
Ask \emph{do you drink tea?} and answer it.
\end{quote}

\subsection*{Your turn: four skills}
\begin{itemize}
\raggedright
  \item \emph{Hear:} \emph{do you go home every day?} Say \textbf{\pa{ਕੀ} \pa{ਤੁਸੀਂ} \pa{ਰੋਜ਼} \pa{ਘਰ} \pa{ਜਾਂਦੇ} \pa{ਹੋ}?}
  \item \emph{Hear:} \emph{no, I don't go home every day} Answer as yourself.
  \item \emph{Read it:} \textbf{\pa{ਕੀ} \pa{ਉਹ} \pa{ਰੋਟੀ} \pa{ਖਾਂਦੀ} \pa{ਹੈ}?} Give its meaning: does she eat roti?
  \item \emph{Write it:} cover the line and write \textbf{\pa{ਕੀ} \pa{ਤੁਸੀਂ} \pa{ਦੁੱਧ} \pa{ਪੀਂਦੇ} \pa{ਹੋ}?}
\end{itemize}

\begin{tcolorbox}[breakable,colback=teal!4,colframe=teal!35!black,title={Before you move on}]
Hold the whole four-line exchange once more, both parts, without the model.
\end{tcolorbox}
